\documentclass{article}
\usepackage{graphicx} % Required for inserting images
\usepackage{tikz}
\usepackage{physics}
\usepackage{gensymb}
\usepackage{amsmath, amsfonts, mathtools, amsthm, amssymb, mathrsfs}
\usepackage{titlesec}
\usepackage{fancyhdr}
\usepackage{hyperref}
\usepackage{float}
\usepackage{booktabs}
\usepackage{enumitem}
\usepackage{subcaption}
\usepackage{multicol}
\usepackage{import}
\usepackage[usenames,dvipsnames]{xcolor}
\usepackage[a4paper, margin=1in]{geometry}
% The configuration for boxes and parts of the overall idea come from Pingbang Hu: https://github.com/sleepymalc
% Use a relative path (forward slashes or "..") to avoid backslash-escapes
\subimport{../}{header.tex}

\title{\textbf{EESC 214 Earth Science Data Analysis}}
\author{Autumn Weaver}
\date{Fall 2026}

\begin{document}

\maketitle

\tableofcontents{\textbf{}}

\pagebreak

\pagestyle{fancy}
\fancyhead{}
\fancyhead[L]{\thetitle}
\fancyhead[R]{Autumn Weaver}

\section{Looking at Data}
Objectives when first looking at data:
\begin{enumerate}
    \item Understand the general character of the dataset 
    \item Understand the general behavior of individual parameters
    \item Detect obvious problems with the data
\end{enumerate}


Several steps when analyzing data:
\begin{enumerate}
    \item Format conversion \\
    Format conversion is often needed but can introduce error. 
    \item Reality Checks 
    
\end{enumerate}



\end{document}